\chapter{結論}
\label{chap:結論}

\section{研究のまとめ}

本研究では，「すたすた」「のろのろ」といったオノマトペを入力として，
ヒューマノイドロボットの歩き方を直感的に切り替える手法を提案した．
手法は 2 段階から構成される．
まず Phase~1 の対照学習でオノマトペと歩容モーションの意味的な対応を獲得し，
続く Phase~2 の強化学習でその対応を物理シミュレーション上の歩行制御に接続した．

Phase~1 では，2D の歩容系列とオノマトペのテキスト埋め込みを
共通の潜在空間に配置する対照学習を行った．
潜在空間上ではクラスタが完全に分離するには至らなかったが，
テキスト埋め込みを基準とした検索では語彙と歩容の対応が一定の水準で成立しており，
言語-動作間の意味的な整合を獲得できることが確認された．
また，学習に含めなかった未知のオノマトペについても，
意味的に近い既知ラベルの近傍に配置される傾向が見られ，
新しい語彙への汎化の可能性が示された．

Phase~2 では，Phase~1 の潜在空間をスタイル報酬として強化学習に組み込み，
オノマトペに応じた歩行生成を試みた．
語彙間で速度やコサイン類似度に差が観測され，
条件付けが歩行に影響を与えていること自体は確認できた．
しかし，語彙ごとに見た目で区別できるほどの歩容差には至っておらず，
未知語を入力した場合は既知語で見られた速度差も再現されなかった．
考察では，どの語彙でもそれなりに高い報酬を得られる
「無難な歩き方」に方策が収束しやすい点を主な原因として指摘した．

\section{今後の課題}

本研究で残された課題と，それに対する今後の方向性を整理する．

まず，語彙間のスタイル分離が不十分である点が最も大きな課題である．
現行の報酬は正例との類似度のみで評価しているため，
語彙間の差を積極的に広げる力が弱い．
考察で述べた hard-negative 型の報酬設計や，
カリキュラム学習による段階的なスタイル圧の導入が有効と考えられる．

次に，2D 歩容系列の情報制約がある．
HOYO は正面視点の 2D データであり，歩幅や奥行き方向の動きを直接扱えない．
この制約がスタイル差の学習を難しくしている面があり，
3D モーションキャプチャデータを教師に用いることで
大きな改善が期待できる．

また，姿勢品質の問題も残っている．
一部の試行で上肢の姿勢が崩れており，
速度に差があっても見た目に歩き方が違うとは感じにくい場合があった．
教師モーションとの関節角誤差を報酬に組み込むなど，
姿勢の質を直接制御する仕組みが必要である．

最後に，未知語への一般化と評価方法の整備がある．
対照学習側では未知語が既知語の近傍に配置される傾向が見られたが，
強化学習側での歩行出力にはその差が十分に反映されていない．
未知語を含む固定の評価セットを用意し，
既知語と同一条件で比較するプロトコルの整備が求められる．
あわせて，定量指標だけでは「歩き方が違って見えるか」を捉えきれないため，
人間による主観評価の導入も重要な課題である．
